\begin{abstract}
 % We study decision making with structured observation (DMSO). The complexity for DMSO has been characterized by a series of work \citep{foster2021statistical, chen2022unified,foster2023tight}. Still, there is a gap between known regret upper and lower bounds: current upper bounds incur a model estimation error that scales with the size of the model class. The work of \cite{foster2024model} made an initial attempt to reduce the estimation error to only scale with the size of the value function set, resulting in the complexity called optimistic decision-estimation coefficient (optimistic DEC). Yet, their approach relies on the \emph{optimism principle} to drive exploration, which deviates from the general idea of DEC that drives exploration only through information gain. 

 %    In this work, we introduce an improved model-free DEC, called \decname, that removes the optimism mechanism in \cite{foster2024model}, making it more aligned with existing model-based DEC. \decname is always upper bounded by optimistic DEC, and could be significantly smaller in special cases. Importantly, the removal of optimism allows it to seamlessly handle \emph{adversarial} environments without explicitly constructing reward estimators, while it was unclear how to achieve it within the optimistic DEC framework. By applying \decname to hybrid MDPs where the transition is stochastic but the reward is adversarial with known features, we provide the first \emph{model-free} regret bounds in \emph{hybrid MDPs} with \emph{bandit feedback} in bilinear classes and Bellman-complete MDPs with bounded coverability, resolving the main open problem in \cite{liu2025decision}.
    
 %    We also improve online function-estimation procedure used in model-free learning: For average estimation error minimization, we improve the estimator to achieve better concentration. This improves the $T^{\frac{3}{4}} $and $T^{\frac{5}{6}}$ regret of \cite{foster2024model} to $T^{\frac{2}{3}} $ and $T^{\frac{7}{9}}$ for the cases with on-policy and off-policy estimation, respectively. For squared estimation error minimization in Bellman-complete MDPs, we redesign the two-timescale procedure in \cite{agarwal2022non} and \cite{foster2024model}, achieving $\sqrt{T}$ regret that improves over the $T^{\frac{2}{3}}$ regret by \cite{foster2024model}. This is also the first time the performance of a DEC-based approach for Bellman-complete MDPs matches that of optimism-based approaches \citep{jin2021bellman, xie2022role}. 



% We study decision making with structured observation (DMSO), whose complexity has been characterized by a series of works \citep{foster2021statistical, chen2022unified, foster2023tight}. Still, a gap remains between known upper and lower regret bounds: current upper bounds depend on model estimation error scaling with the size of the model class. \citet{foster2024model} attempted to reduce this dependence to the value-function class via the optimistic decision-estimation coefficient (optimistic DEC), but their approach relies on optimism-based exploration rather than pure information gain.

% We introduce a new model-free coefficient, \decname, which removes optimism from \citet{foster2024model} and aligns with the general DEC framework. \decname\ is always upper bounded by optimistic DEC and can be substantially smaller in special cases. Crucially, eliminating optimism enables direct handling of adversarial rewards without explicit reward estimators. Applying \decname\ to hybrid MDPs with stochastic transitions and adversarial rewards, we obtain the first model-free regret bounds for hybrid MDPs with bandit feedback under bilinear and Bellman-complete settings, resolving the main open problem in \citet{liu2025decision}.

% We further refine online estimation in model-free learning. For average estimation error minimization, our improved estimator enhances concentration, tightening regret from $T^{3/4}$ and $T^{5/6}$ to $T^{2/3}$ and $T^{7/9}$ in on-policy and off-policy cases, respectively. For squared estimation error minimization in Bellman-complete MDPs, our redesigned two-timescale update achieves $\sqrt{T}$ regret—matching optimism-based approaches \citep{jin2021bellman, xie2022role} for the first time within the DEC framework.

%We study decision making with structured observation (DMSO), whose complexity decision-estimation coefficient (DEC) has been characterized by previous works \citep{foster2023tight}. Still, a gap remains between known upper and lower regret bounds: existing upper bounds depend on estimation errors scaling with the size of the model class. \citet{foster2024model} attempted to reduce this dependence to the value-function class via the optimistic DEC, but their approach relies on optimism-based exploration which can only handle stochastic settings.

We study decision making with structured observation (DMSO). Previous work \citep{foster2021statistical, foster2023tight} has characterized the complexity of DMSO via the decision-estimation coefficient (DEC), but left a gap between the regret upper and lower bounds that scales with the size of the model class. To tighten this gap, \cite{foster2024model} introduced optimistic DEC, achieving a bound that scales only with the size of the value-function class. However, their optimism-based exploration is only known to handle the stochastic setting, and it remains unclear whether it extends to the adversarial setting.

%To tackle their limitations, we propose \decname, a model-free DEC that removes optimism and restores the pure information-gain principle of DEC. \decname\ is always no larger than optimistic DEC and can be much smaller in special cases. Its optimism-free nature enables handling adversarial rewards without explicit reward estimation. Applied to hybrid MDPs with stochastic transitions and adversarial rewards, \decname\ yields the first model-free regret bounds for hybrid MDPs with bandit feedback and general function approximation, resolving the main open problem of \citet{liu2025decision}.

We introduce \decname, a model-free DEC that removes optimism and drives exploration purely by information gain. \decname is always no larger than optimistic DEC and can be much smaller in special cases. Importantly, the removal of optimism allows it to handle adversarial environments without explicit reward estimators. By applying \decname to hybrid MDPs with stochastic transitions and adversarial rewards, we obtain the first \emph{model-free} regret bounds for \emph{hybrid} MDPs with \emph{bandit} feedback under linear reward and several \emph{general} transition structures, resolving the main open problem left by \cite{liu2025decision}.

%We also improve online function-estimation procedure used in model-free learning: For average estimation  error minimization, we refine the estimator to achieve sharper concentration, improving the regret bounds of \cite{foster2024model} from $T^{\frac{3}{4}}$ and $T^{\frac{5}{6}}$ to $T^{\frac{2}{3}}$ and $T^{\frac{7}{9}}$ for on-policy and off-policy estimation, respectively. For squared error minimization in Bellman-complete MDPs, we redesign the two-timescale scheme of \cite{agarwal2022non} and \cite{foster2024model}, attaining $\sqrt{T}$ regret, surpassing the $T^{\frac{2}{3}}$ rate in \cite{foster2024model} and marking the first DEC-based method for Bellman-complete MDPs matching optimism-based results \citep{jin2021bellman, xie2022role}.

We also improve the online function-estimation procedure in model-free learning: For average estimation error minimization, we refine \cite{foster2024model}'s estimator to achieve sharper concentration, improving their regret bounds from $T^{\frac{3}{4}}$ to $T^{\frac{2}{3}}$ (on-policy) and from $T^{\frac{5}{6}}$ to $T^{\frac{7}{9}}$ (off-policy). For squared error minimization in Bellman-complete MDPs, we redesign their two-timescale procedure, improving the regret bound from $T^{\frac{2}{3}}$ to $\sqrt{T}$. This is the first time a DEC-based method achieves performance matching that of optimism-based approaches \citep{jin2021bellman, xie2022role} in Bellman-complete MDPs. 



% We study decision making with structured observation (DMSO), whose complexity decision-estimation coefficient (DEC) has been characterized by previous works (Foster et al., 2023a). Still, a gap remains between known upper and lower regret bounds: existing upper bounds depend on estimation errors scaling with the size of the model class. Foster et al. (2023b) attempted to reduce this dependence to the value-function class via the optimistic DEC, but their approach relies on optimism-based exploration which can only handle stochastic settings.

% To tackle their limitations, we propose Dig-DEC, a model-free DEC that removes optimism and restores the pure information-gain principle of DEC. Dig-DEC is always no larger than optimistic DEC and can be much smaller in special cases. Its optimism-free nature enables handling adversarial rewards without explicit reward estimation. Applied to hybrid MDPs with stochastic transitions and adversarial rewards, Dig-DEC yields the first model-free regret bounds for hybrid MDPs with bandit feedback and general function approximation, resolving the main open problem of Liu et al. (2025).

% We also improve online function-estimation procedure used in model-free learning. For average error minimization, we refine the estimator to achieve sharper concentration, improving the regret bounds of Foster et al. (2023b) from T^{\frac{3}{4}} and T^{\frac{5}{6}} to T^{\frac{2}{3}} and T^{\frac{7}{9}} for on- and off-policy settings, respectively. For squared error minimization in Bellman-complete MDPs, we redesign the two-timescale scheme of Agarwal and Zhang (2022) and Foster et al. (2023b), attaining \sqrt{T} regret, surpassing the T^{\frac{2}{3}} rate in Foster et al. (2023b) and marking the first DEC-based method matching optimism-based results (Jin et al., 2021; Xie et al., 2023).
\end{abstract}